% SOURCE_AUTHORITY: sga5_fr_workpass.tex, lines 12896--12917 (LF-normalized unit).
% SOURCE_SHA256: 791F4EFFC5E02832D5D77ED03518C8156D6F07E4C8238B03545DB93D883FBB28
\subsubsection*{6.2}

Usaremos as seguintes notações:
\[
\alpha_x^\Delta(F)
=
\clK{\Delta}\bigl(\Hom^\bullet_{\Lambda[G]}(Sw_x^\Lambda,F_{\bar\eta})\bigr),
\tag{6.2.1}
\]
\[
\varepsilon_x^\Delta(F)
=
\alpha_x^\Delta(F)+
\bigl(\clK{\Delta}(F_{\bar\eta})-\clK{\Delta}(F_x)\bigr).
\tag{6.2.2}
\]
Naturalmente, na fórmula (6.2.1), $F_{\bar\eta}$ está dotado da estrutura de $A[G]$-módulo mediante o isomorfismo (6.1.1). Observe-se que, se $x\in C(k)$ é tal que $C'$ é moderadamente ramificado em $x$, tem-se (4.6) $Sw_x=0$, logo $\alpha_x^\Delta(F)=0$ e, portanto, neste caso tem-se (de acordo com Ogg--Chafarevitch):
\[
\varepsilon_x^\Delta(F)=\clK{\Delta}(F_{\bar\eta})-\clK{\Delta}(F_x).
\tag{6.2.3}
\]
É claro que $\varepsilon_x^\Delta(F)=0$ se $x\notin Y=C-U$ e, portanto, só existe um número finito de $x$ tais que $\varepsilon_x^\Delta(F)\ne0$.
